% answers_short.tex —— 2005 年·普通高等学校招生考试·天津卷·理科 · 简短答案速查
% 由 tools/make_short_answers.py 自 answers.tex 的 \gksolhead/\gkanswer 提取，
% 并与 2005/2005_数学真题_天津理科_答案要点.md 逐题比对后生成，勿手改（改源 MD 后重跑）。

\begin{gknotice}
\textbf{使用说明}：本节只给最终答案便于快速核对，完整推导见解析版。答案为编写者独立推导，\textbf{非官方标准答案}。
\end{gknotice}

\gkpart{选择题}
\begin{center}\begin{tabular}{ccccccccccc}
\toprule
题号 & 1 & 2 & 3 & 4 & 5 & 6 & 7 & 8 & 9 & 10 \\
\midrule
答案 & D & C & B & D & C & B & A & C & A & B \\
\bottomrule
\end{tabular}\end{center}

\gkpart{填空题}
\begin{itemize}[leftmargin=2.6em]
  \item[\textbf{11.}] $\dfrac{7^n - 1}{6}$。
  \item[\textbf{12.}] $\sqrt{2}$。
  \item[\textbf{13.}] $2600$。
  \item[\textbf{14.}] $\left(-\dfrac{\sqrt{10}}{5}, \dfrac{3\sqrt{10}}{5}\right)$。
  \item[\textbf{15.}] $4760$。
  \item[\textbf{16.}] $0$。
\end{itemize}

\gkpart{解答题}
\begin{itemize}[leftmargin=2.6em]
  \item[\textbf{17.}] $\angle A = \dfrac{\pi}{3}$（或 $60^\circ$），$\tan B = \dfrac{1}{2}$。
  \item[\textbf{18.}] (1) 当 $a = 1$ 时，$S_n = \dfrac{n(n+3)}{2}$；当 $a \neq 1$ 时，$S_n$ 表达式见解析；(2) $\max(a, b)$（当 $a \geqslant b$ 时为 $a$；当 $a < b$ 时为 $b$）
  \item[\textbf{19.}] (1) $A_1A$ 与底面 $ABC$ 所成的角为 $60^\circ$（或 $\dfrac{\pi}{3}$）；(2) 见解析版；(3) $\dfrac{4\sqrt{3}}{27}\pi a^3$
  \item[\textbf{20.}] 此人距水平地面 $60$ 米时，观看塔的视角 $\angle BPC$ 最大。
  \item[\textbf{21.}] (1) 焦点坐标为 $\left(0, \dfrac{1}{4a}\right)$，准线方程为 $y = -\dfrac{1}{4a}$；(2) 见解析版；(3) $y_1$ 的取值范围为 $(-\infty, -1) \cup \left(-1, -\dfrac{1}{4}\right)$
  \item[\textbf{22.}] (1) 见解析版；(2) 见解析版；(3) 见解析版
\end{itemize}
